\documentclass[10pt,a4paper]{amsart}
\usepackage[margin=19mm]{geometry}
\usepackage{amsmath,amssymb,mathtools}
\usepackage[scr=rsfs,cal=euler]{mathalfa}
\usepackage{xfrac}
\usepackage[unicode,hidelinks,bookmarks=false]{hyperref}
\newcommand{\ie}{i.e.\ }
\newcommand{\cf}{cf.\ }
\newcommand{\resp}{respectively\ }
\newcommand{\Subsection}{\subsection}
\providecommand{\nobreakdash}{\nobreak\hbox{-}}
\setlength{\emergencystretch}{2em}
\multlinegap0pt
\usepackage{amssymb}
\usepackage{enumerate}
\usepackage[shortlabels]{enumitem}
\usepackage{cancel}
\newtheorem{corollary}{Corollary}
\newtheorem{lemma}{Lemma}
\newcommand{\bb}{\mathcal{B}}
\newcommand{\eps}{\epsilon}
\newcommand{\ol}{\overline}
\newcommand{\sll}{\mathfrak{sl}}
\newcommand{\sm}{\setminus}
\newcommand{\uvsln}{U_v(\sll_n)}
\title{Representation theory of mirabolic quantum $\mathfrak{sl}_n$}
\author{Pallav Goyal, Daniele Rosso}
\date{Source: arXiv:2510.07469v2 (CC BY 4.0)}
\begin{document}
\maketitle

\noindent\emph{Source and licence.} Representation theory of mirabolic quantum $\mathfrak{sl}_n$. Version arXiv:2510.07469v2 (CC BY 4.0), distributed by arXiv under the Creative Commons Attribution 4.0 International licence, http://creativecommons.org/licenses/by/4.0/, as recorded in the arXiv licence metadata for this exact version and shown on the abstract page https://arxiv.org/abs/2510.07469v2. Full licence notice: \texttt{licenses/arxiv-2510.07469-CC-BY-4.0.txt}.\par\smallskip

\noindent\textit{Editorial scope.} Counted unit: the article's lemma lem:I\_r+1 (source0.tex lines 1098--1102). The article's complete original proof of it is delivered with it (source0.tex lines 1104--1360, one \texttt{proof} environment of 257 lines). No mathematics is written, repaired or re-derived: the statement and the proof are the article's own lines, in the article's own notation, and no retained line is substituted or re-flowed.  Retained uncounted as the source-local context the proof needs: lines 201--206; lines 223--227; lines 327--337; lines 864--872; lines 1046--1047; lines 1050--1052; lines 1056--1057; lines 1061--1061; lines 1065--1065; lines 1068--1071.  Nothing was omitted from any retained span, so the package carries no omission record.  No retained line contains a citation, so the wrapper carries no bibliography and no bibliography entry.  No retained line contains a figure environment, an image-inclusion command or drawing code, so no image is delivered and the package declares no asset dependency.  Editorial formatting only, in addition: an AMS \texttt{amsart} preamble that restates the article's own declarations and macros used by the retained lines, namely the environments \texttt{corollary}, \texttt{lemma} on separate counters, together with the standard packages those lines need.  The retained environments print in this document as Corollary 1, Lemma 1 in order of appearance, whereas the article prints the same items with the numbering of its own full document.  The complete native source of the article as distributed in the arXiv e-print for this exact version is bundled unchanged as \texttt{sources/186672/source0.tex}.
\begin{align}\label{eq:usln.1} k_i^{\pm1}k_i^{\mp1}=1, & \quad k_ik_j=k_jk_i \\
\label{eq:usln.2} k_ie_jk_i^{-1} = v^{a_{ij}}e_j, & \quad k_if_jk_i^{-1} = v^{-a_{ij}}f_j\\
\label{eq:usln.3} e_if_j-f_je_i & =  \delta_{ij}\frac{k_i-k_i^{-1}}{v-v^{-1}}\\
\label{eq:usln.4} e_ie_j=e_je_i, & \quad f_if_j=f_jf_i \quad\text{ if }|i-j| \neq 1 \\
\label{eq:usln.5} (v+v^{-1})e_ie_je_i & =  e_i^2e_j + e_je_i^2 \quad\text{ if }|i-j| = 1\\
\label{eq:usln.6} (v+v^{-1})f_if_jf_i & =  f_i^2f_j + f_jf_i^2 \quad\text{ if }|i-j| = 1.\end{align}

\begin{align}\label{eq:mun.1} \ell^2=\ell,&\quad k_i\ell=\ell k_i\\
\label{eq:mun.2} \ell e_i=e_i\ell,&\quad \ell f_i=f_i\ell \quad\text{ if }i\geq 2\\
\label{eq:mun.3} \ell e_1\ell=\ell e_1,&\quad \ell f_1\ell=f_1\ell\\
\label{eq:mun.4} (v+v^{-1})e_1\ell e_1 &= v^{-1}e_1^2\ell + v\ell e_1^2\\
\label{eq:mun.5} (v+v^{-1})f_1\ell f_1 &= v^{-1}\ell f_1^2 + vf_1^2\ell.\end{align}

\begin{corollary} \label{cor:lastcor}
Consider any non-commutative monomial $Q(f)\in \uvsln$ in the $f_i$'s.
\begin{enumerate}
\item The element $Q(f)$ can be expressed a linear combination of monomials (having the same length as $Q(f)$) of the form:
\[f_1f_2\cdots f_s,\qquad f_1^2 Q'(f), \qquad f_t Q'(f)\]
for some $0\leq s\leq n-1$, $1<t\leq n-1$ and some non-commutative monomials $Q'(f)$.
\item The element $\ell Q(f)\in MU(n)$ can be expressed a linear combination of monomials (having the same length as $\ell Q(f)$) of the form:
\[Q'(f)\ell f_1f_2\cdots f_s\]
for some $0\leq s\leq n-1$ and some non-commutative monomials $Q'(f)$.
\end{enumerate}
\end{corollary}

\begin{corollary} \label{cor:fmove}
Let $N$ be an arbitrary $MU(n)$-module and let $y\in N$ be a mirabolic highest weight vector with depth $r$. Consider the vector:
\[y':= f_{i_1} f_{i_2}\cdots f_{i_k}  \ell f_1 f_2\cdots f_{s-1} f_{s} \cdot y,\]
for some $i_j$ such that $1\leq i_1< i_2<\cdots <i_k\leq r$ and some $s\geq r$. Then, we have the equality:
\[y'=  \frac{c}{(v+v^{-1})^k}\ell  f_1^{\eps_1}  f_2^{\eps_2} \cdots f_s^{\eps_s} \cdot y,\]
where:
\[\eps_t=\begin{cases}2 & \text{if}~t=i_j~\text{for some}~j,\\1 &\text{otherwise},\end{cases}\]
and $c=v^{-1}$ if $i_1=1$ and $c=1$ otherwise.
\end{corollary}

\begin{equation}\label{eq:Phi1}\Phi(1,I_1) := v^{\lambda_1-\lambda_{r+1}}\ell f_1f_2\cdots f_r \cdot  w_{\lambda,r};\qquad
\Phi(1,I_i):=e_{i-1}e_{i-2}\cdots e_1\cdot \Phi(1,I_1), \text{ for }2\leq i\leq r+1.\end{equation}%Then, for any $2\leq i\leq r+1$, define $\Phi(1,I_i):=e_{i-1}e_{i-2}\cdots e_1\cdot \Phi(1,I_1)$. 

\begin{align}\notag\Phi(x_{\lambda,r})=\Phi(1,\{1,\ldots,r\})& =\Phi(1,I_{r+1})=v^{\lambda_1-\lambda_{r+1}}e_re_{r-1}\cdots e_1\ell f_1\cdots f_{r-1}f_r\cdot w_{\lambda, r}\\
\label{eq:topvect}\text{ (by the definition of $V_{\lambda,r}$)} & = w_{\lambda,r}.
\end{align}

\begin{align}\label{eq:Phi-fj}\Phi(f_jQ(f),I_1)&:=f_j\cdot \Phi(Q(f),I_1),~\text{ for all }j\neq r+1; \\
\label{eq:Phi-fr}\Phi(f_{r+1}Q(f),I_{1})&:=v^{\mu_1-\mu_{r+1}+1}\ell f_1f_2\cdots f_{r+1} \cdot \Phi(Q(f),I_{r+1}).\end{align}

\begin{equation}\label{eq:Phi-I+1}\Phi(\widetilde Q(f),I_{i+1})=e_i\cdot \Phi(\widetilde Q(f), I_i)-v^{-1}\Phi([e_i,\widetilde Q(f)], I_{i}).\end{equation}

\begin{equation}\label{eq:Phi-I'}\Phi(Q(f),I):=v^{\mu_{i_j-1}-\mu_{i_j}}f_{i_j-1}\cdot \Phi(Q(f),I') - v^{\mu_{i_j-1}-\mu_{{i_j}}}\Phi(f_{i_j-1}Q(f), I').\end{equation}

\begin{lemma}\label{lem:Phi-ei}
For all $1\leq i\leq r$ and $Q(f)\in \bb$ we have the equality:
\[e_i\cdot \Phi(Q(f), (I_{i}\cup\{r+2\})\sm \{r+1\}) =  \Phi(e_i\cdot(Q(f), (I_{i}\cup\{r+2\})\sm \{r+1\})).\]
\end{lemma}

\begin{lemma} \label{lem:I_r+1}
For any $1\leq i\leq r+1$ and any $Q(f)\in \mathcal B$, we have the equality
\[\Phi(z\cdot(Q(f), I_i)) = z\cdot \Phi(Q(f), I_i),\] where $z=\ell,e_j,f_j$ %or $z=f_j$ 
for $1\leq j\leq n-1$.
\end{lemma}

\begin{proof}
We will prove the lemma for pairs $(Q(f), I_i)$ via induction on $\ol\deg(Q(f))+i$. Throughout the rest of the proof, we will assume that $Q(f)$ has weight $\mu-\lambda$ with respect to the $k_j$'s. We start with the base case of the induction when $i=1$ and $Q(f)=1$. When $z=f_j$ and $j\neq r+1$, we have that:
\begin{equation}\label{eq:fj-11}\Phi(f_j\cdot (1, I_1)) = \Phi(f_j, I_1) = f_j\cdot \Phi(1,I_{1}),\end{equation}
where the first equality follows from the definition of $\rho$ and second equality is by \eqref{eq:Phi-fj}. %the definition of $\Phi(f_jQ(f), I_1)$. 
When $j=r+1$, 
\begin{align}
\notag\Phi(f_{r+1}\cdot (1, I_1)) &= \Phi(f_{r+1}, I_1) + v^{\lambda_{r+2}-\lambda_{r+1}}\Phi(1,(I_1\cup\{r+2\})\sm\{r+1\})\\
\label{eq:fr-11}\text{(by \eqref{eq:Phi-I'})}&= f_{r+1}\cdot \Phi(1,I_{1}).
\end{align}
Next, we suppose $z=e_j$ for some $j$. If $j=1$, we have:
\[\Phi(e_1\cdot (1,I_1))  = \Phi(1,I_2)\\
 = e_1\cdot \Phi(1,I_1),\]
where the second equality is by \eqref{eq:Phi1}. Now, we suppose that $j\neq 1$. Note that $e_j\cdot (1,I_1)=0$. On the other hand, if $j>r$ we have
\[e_j\cdot\Phi(1,I_1)= v^{\lambda_1-\lambda_{r+1}}e_j\ell f_1f_2\cdots f_r\cdot w_{\lambda,r}=v^{\lambda_1-\lambda_{r+1}}\ell f_1f_2\cdots f_re_j\cdot w_{\lambda,r}=v^{\lambda_1-\lambda_{r+1}}\ell f_1f_2\cdots f_r\cdot( e_j\cdot w_{\lambda,r})=0;\]
and, when $2\leq j\leq r$, we also have that
\begin{align*}
e_j\cdot\Phi(1,I_1)&= v^{\lambda_1-\lambda_{r+1}}e_j\ell f_1f_2\cdots f_r\cdot w_{\lambda,r}=v^{\lambda_1-\lambda_{r+1}}\ell e_j f_1f_2\cdots f_r\cdot w_{\lambda,r}\\
&=v^{\lambda_1-\lambda_{r+1}}\ell f_1f_2\cdots f_re_j\cdot w_{\lambda,r} + v^{\lambda_1-\lambda_{r+1}}\ell f_1f_2\cdots f_{j-1}\frac{k_j-k_j^{-1}}{v-v^{-1}}f_{j+1\cdots}f_r\cdot w_{\lambda,r}\\
&=0,
\end{align*}
where the last equality follows since $e_j$ and $\ell f_1f_2\cdots f_{j-1}$ act on $w_{\lambda, r}$ by zero. Finally, we note that $\ell$ acts on $\Phi(1,I_1) = v^{\lambda_1-\lambda_{r+1}}\ell f_1f_2\cdots f_r \cdot w_{\lambda, r}$ by $1$, proving equivariance with respect to $\ell$ and completing the proof of the base case of the induction.

Next we proceed with the induction step. First, we suppose that $i=1$. The fact that $f_j\cdot\Phi(Q(f),I_1)=\Phi(f_j\cdot(Q(f),I_1))$ follows from exactly the same calculation as when $Q(f)=1$ in \eqref{eq:fj-11} and \eqref{eq:fr-11}. Next, we suppose $z=e_j$ for some $j$. If $j=1$, we have:
\begin{align*}\Phi(e_1\cdot (Q(f),I_1)) & = \Phi(Q(f),I_2) + v^{-1}\Phi(e_1Q(f),I_1) \\
&= \Phi(Q(f),I_2) + v^{-1}\Phi([e_1,Q(f)],I_1) \\
\text{(by \eqref{eq:Phi-I+1})} &= e_1\cdot \Phi(Q(f),I_1).\end{align*}
%where the second equality is by the definition of $\Phi(Q(f), I_2)$.
Now, we suppose that $j\neq 1$. If $\ol\deg(Q(f))\geq 1$, then $Q(f)$ can be written as a linear combination of elements of the form $f_{j'}\widetilde Q(f)$ for varying $j'$ and $\ol\deg(\widetilde Q(f))< \ol\deg(Q(f))$. In this case, if $j'\neq r+1$, we have
\begin{align*}
\Phi(e_j\cdot(f_{j'}\widetilde Q(f), I_1))&=\Phi((e_j\otimes k_j)\cdot(f_{j'}\widetilde Q(f),I_1)+(1\otimes e_j)\cdot(f_{j'}\widetilde Q(f),I_1))\\
\text{(because $e_j\cdot u_{I_1}=0$)} &=v^{\delta_{j,r+1}}\Phi([e_j,f_{j'}\widetilde Q(f)], I_1)\\
&=\delta_{j,j'}\Phi\Big(\frac{k_{j}-k_{j}^{-1}}{v-v^{-1}}\widetilde Q(f),I_1\Big) + v^{\delta_{j,r+1}}\Phi(f_{j'}[e_j,\widetilde{Q}(f)],I_1)\\
\text{(by \eqref{eq:Phi-fj})}&=\delta_{j,j'}\frac{k_{j}-k_{j}^{-1}}{v-v^{-1}}\cdot\Phi(\widetilde Q(f),I_1) + v^{\delta_{j,r+1}}f_{j
'}\cdot \Phi([e_j,\widetilde{Q}(f)],I_1)\\
\text{(by induction and because $e_j\cdot u_{I_1}=0$)}&=\delta_{j,j'}\frac{k_{j}-k_{j}^{-1}}{v-v^{-1}}\cdot\Phi(\widetilde Q(f),I_1) + f_{j'}e_j\cdot\Phi(\widetilde{Q}(f),I_1)\\
&=e_jf_{j'}\cdot \Phi(\widetilde{Q}(f),I_1)\\
\text{(by \eqref{eq:Phi-fj})}&=e_j\cdot \Phi(f_{j'}\widetilde{Q}(f),I_1),
\end{align*}
proving the claim. Finally, suppose $j'=r+1$. In that case:
\begin{align*}
\Phi(e_j\cdot(f_{r+1}\widetilde Q(f), I_1))&=v^{\delta_{j,r+1}}\Phi([e_j,f_{r+1}\widetilde Q(f)], I_1)\\
&=\delta_{j,r+1}v\Phi\Big(\frac{k_j-k_j^{-1}}{v-v^{-1}}\widetilde Q(f),I_1\Big) + v^{\delta_{j,r+1}}\Phi(f_{r+1}[e_j,\widetilde{Q}(f)],I_1)\\
\text{(by \eqref{eq:Phi-fr})}&=\delta_{j,r+1}v\Phi\Big(\frac{k_j-k_j^{-1}}{v-v^{-1}}\widetilde Q(f),I_1\Big) \\
&~+ v^{\delta_{j,r+1}+(\mu_1)-(\mu_{r+1}-\delta_{j,r}+\delta_{j,r+1}+1)+1}\ell f_1f_2\cdots f_rf_{r+1}\cdot \Phi([e_j,\widetilde{Q}(f)],I_{r+1})\\
&=\delta_{j,r+1}v\Phi\Big(\frac{k_j-k_j^{-1}}{v-v^{-1}}\widetilde Q(f),I_1\Big) + v^{\mu_1-\mu_{r+1}}\ell f_1f_2\cdots f_rf_{r+1}e_j\cdot\Phi(\widetilde{Q}(f),I_{r+1}),
\end{align*}
where the last equality follows by observing that $k_j\cdot u_{I_{r+1}}=v^{\delta_{j,r}}u_{I_{r+1}}.$ If $j\neq r+1$, the term on the left is zero, whereas the term on the right becomes:
\begin{align*}
v^{\mu_1-\mu_{r+1}}\ell f_1f_2\cdots f_rf_{r+1}e_j\cdot\Phi(\widetilde{Q}(f),I_{r+1})&=v^{\mu_1-\mu_{r+1}}\ell f_1f_2\cdots f_je_j\cdots f_rf_{r+1}\cdot\Phi(\widetilde{Q}(f),I_{r+1})\\
&=v^{\mu_1-\mu_{r+1}}\ell f_1f_2\cdots \left(e_jf_j -\frac{k_j-k_j^{-1}}{v-v^{-1}}\right)\cdots f_rf_{r+1}\cdot\Phi(\widetilde{Q}(f),I_{r+1})\\
%&\text{(Here, we use the fact that }\Phi(\widetilde{Q}(f),I_{r+1}) \text{ has depth } r \text{ in order to skip the commutator term}.)\\
&=v^{\mu_1-\mu_{r+1}}\ell f_1f_2\cdots e_jf_j \cdots f_rf_{r+1}\cdot\Phi(\widetilde{Q}(f),I_{r+1})\\
%\text{(because $\Phi(\widetilde{Q}(f),I_{r+1})$ has depth $r$ } & \text{by induction on the second part of the Lemma)} \\
&\text{(because } \ell f_1f_2\cdots f_{j-1}\cdot \Phi(\widetilde{Q}(f),I_{r+1})=0 \text{ by the induction step)}\\
&=v^{\mu_1-\mu_{r+1}}e_j\ell f_1f_2\cdots f_rf_{r+1}\cdot\Phi(\widetilde{Q}(f),I_{r+1})\\
\text{(by \eqref{eq:Phi-fr})}&=e_j\cdot \Phi(f_{r+1} \widetilde{Q}(f),I_{1}).
\end{align*}
Finally, if $j=r+1$, the sum equals:
\begin{align*}
&~v\Phi\Big(\frac{k_{r+1}-k_{r+1}^{-1}}{v-v^{-1}}\widetilde Q(f),I_1\Big) + v^{\mu_1-\mu_{r+1}}\ell f_1f_2\cdots f_rf_{r+1}e_{r+1}\cdot\Phi(\widetilde{Q}(f),I_{r+1})\\
=&~v\Phi\Big(\frac{k_{r+1}-k_{r+1}^{-1}}{v-v^{-1}}\widetilde Q(f),I_1\Big) + v^{\mu_1-\mu_{r+1}}e_{r+1}\ell f_1f_2\cdots f_rf_{r+1}\cdot\Phi(\widetilde{Q}(f),I_{r+1})\\
&-v^{\mu_1-\mu_{r+1}}\ell f_1f_2\cdots f_r \cdot\Phi\Big(\frac{k_{r+1}-k_{r+1}^{-1}}{v-v^{-1}}\widetilde{Q}(f),I_{r+1}\Big)\\
=&~v\Phi\Big(\frac{k_{r+1}-k_{r+1}^{-1}}{v-v^{-1}}\widetilde Q(f),I_1\Big) + v^{\mu_1-\mu_{r+1}}e_{r+1}\ell f_1f_2\cdots f_rf_{r+1}\cdot\Phi(\widetilde{Q}(f),I_{r+1}) -  v\Phi\Big(\frac{k_{r+1}-k_{r+1}^{-1}}{v-v^{-1}}\widetilde{Q}(f),I_{1}\Big)\\
&(\text{where the last equality follows from the fact that }\\
&\ell f_1f_2\cdots f_r \cdot\Phi\Big(\frac{k_{r+1}-k_{r+1}^{-1}}{v-v^{-1}}\widetilde{Q}(f),I_{r+1}\Big) = \Phi\Big(\ell f_1f_2\cdots f_r \cdot\left(\frac{k_{r+1}-k_{r+1}^{-1}}{v-v^{-1}}\widetilde{Q}(f),I_{r+1}\right)\Big)\\
&\text{by the induction step)}\\
=&~v^{\mu_1-\mu_{r+1}}e_{r+1}\ell f_1f_2\cdots f_rf_{r+1}\cdot\Phi(f_{r+1}\widetilde{Q}(f),I_{r+1})\\
\text{(by \eqref{eq:Phi-fr})}=&~e_{r+1}\cdot \Phi(f_{r+1}\widetilde{Q}(f),I_1).
\end{align*}

Lastly, we suppose $z=\ell$. We need to show that $\ell$ acts on $\Phi(Q(f), I_1)$ by $1$. We write $Q(f)$ as a linear combination of elements of the form $f_j\widetilde{Q}(f)$,with $\ol{\deg}(\widetilde{Q}(f))<\ol{\deg}(Q(f))$. If $j\neq r+1$, we have, by \eqref{eq:Phi-fj},
\[\ell\cdot \Phi(f_j\widetilde Q(f),I_1) =\ell f_j\cdot \Phi(\widetilde Q(f),I_1)= \ell f_j\ell\cdot \Phi(\widetilde Q(f),I_1) = f_j\ell\cdot \Phi(\widetilde Q(f),I_1)=f_j\cdot \Phi(\widetilde Q(f),I_1)=\Phi(f_j\widetilde Q(f),I_1)\]
where the second and fourth equality follow by the induction step.
Finally, by \eqref{eq:Phi-fr},
\begin{align*}
\ell\cdot \Phi(f_{r+1}\widetilde Q(f),I_1) &= \ell\cdot(v^{\mu_1-\mu_{r+1}+1}\ell f_1f_2\cdots f_{r+1} \cdot \Phi(\widetilde Q(f),I_{r+1}))\\
&= v^{\mu_1-\mu_{r+1}+1}\ell f_1f_2\cdots f_{r+1} \cdot \Phi(\widetilde Q(f),I_{r+1})\\
&= \Phi(f_{r+1}\widetilde Q(f),I_1),
\end{align*}completing the proof.

Now, we work with general $i$. Recall that, by \eqref{eq:Phi-I+1}, 
\[\Phi( Q(f),I_{i+1})=e_i\cdot \Phi( Q(f), I_i)-v^{-1}\Phi([e_i,Q(f)], I_{i}).\]

We need to show that $\Phi(z\cdot(Q(f), I_{i+1})) = z\cdot \Phi(Q(f), I_{i+1})$ where $z$ is a generator of $MU(n)$. We accomplish this in $3$ steps, dealing with the $f_j$'s, $\ell$ and $e_j$'s separately.

\underline{\textbf{Step $1$: $z=f_j$}}

Let $z=f_j$ and $j\not\in\{i-1,i,r+1\}$. Then,
\begin{align*}\Phi( f_j\cdot(Q(f),I_{i+1}))=\Phi( f_jQ(f),I_{i+1})&=e_i\cdot \Phi( f_jQ(f), I_i)-v^{-1}\Phi([e_i,f_jQ(f)], I_{i})\\
\text{(because $f_j\cdot u_{I_i}=0$ and $[e_i,f_j]=0$) }&=e_i\cdot \Phi(f_j\cdot(Q(f), I_i))-v^{-1}\Phi(f_j[e_i,Q(f)], I_{i})\\
\text{(by the induction step and because $f_j\cdot u_{I_i}=0$) }&=e_if_j\cdot \Phi(Q(f), I_i)-v^{-1}\Phi(f_j\cdot([e_i,Q(f)], I_{i}))\\
\text{(by the induction step and because $[e_i,f_j]=0$) }&=f_je_i\cdot \Phi(Q(f), I_i)-v^{-1}f_j\cdot\Phi([e_i,Q(f)], I_{i}) \\
&=  f_j\cdot\Phi(Q(f),I_{i+1}).\end{align*}
%we get the required claim by acting by $f_j$ on both sides of the above equality. 
Next, we first consider the case when $i\geq 2$ and $z=f_{i-1}$. Then,
\begin{align*}
f_{i-1}\cdot\Phi( Q(f),I_{i+1})&=e_if_{i-1}\cdot \Phi( Q(f), I_i)-v^{-1}f_{i-1}\cdot\Phi([e_i,Q(f)], I_{i})\\
\text{(by induction hypothesis) }&=e_i\cdot \Phi(f_{i-1}\cdot( Q(f), I_i))-v^{-1}\Phi(f_{i-1}\cdot([e_i,Q(f)], I_{i}))\\
&=e_i\cdot \Phi(f_{i-1}Q(f),I_i) + v^{\mu_{i}-\mu_{i-1}}e_i\cdot \Phi(Q(f), I_{i-1})\\
&~-v^{-1}\Phi(f_{i-1}[e_i,Q(f)], I_{i})-v^{\mu_i-\mu_{i-1}}\Phi([e_i,Q(f)], I_{i-1})\\
\text{(by induction hypothesis and because $e_i\cdot u_{I_{i-1}}=0$) }&=\Phi(f_{i-1}Q(f), I_{i+1})\\
\text{(because $f_{i-1}\cdot u_{I_{i+1}}=0$) } &=\Phi(f_{i-1}\cdot(Q(f), I_{i+1})).
\end{align*}
Next, when $z=f_i$, we get:
\begin{align*}
f_{i}\cdot\Phi( Q(f),I_{i+1})&=f_ie_i\cdot \Phi( Q(f), I_i)-v^{-1}f_{i}\cdot\Phi([e_i,Q(f)], I_{i})\\
\text{(by induction hypothesis and because $f_i\cdot u_{I_i}=0$) }&=e_if_i\cdot \Phi( Q(f), I_i)-\frac{k_i-k_i^{-1}}{v-v^{-1}}\cdot \Phi( Q(f), I_i)-v^{-1}\Phi(f_i[e_i,Q(f)], I_{i})\\
\text{(by induction hypothesis and because $f_i\cdot u_{I_i}=0$) }&=e_i\cdot \Phi(f_iQ(f), I_i)-\frac{k_i-k_i^{-1}}{v-v^{-1}}\cdot \Phi( Q(f), I_i)\\
&~-v^{-1}\Phi([e_i,f_iQ(f)], I_{i})+v^{-1}\Phi\Big(\frac{k_i-k_i^{-1}}{v-v^{-1}}Q(f), I_{i}\Big)\\
\text{(because $\rho(k_i^{\pm 1})=k_i^{\pm 1}\otimes k_i^{\pm 1}$)}&=e_i\cdot \Phi(f_iQ(f),I_i)-v^{-1}\Phi([e_i,f_iQ(f)],I_i)+v^{-1}k_i^{-1}\cdot \Phi(Q(f),I_i)\\
&=\Phi(f_iQ(f),I_{i+1})+ \Phi(k_i^{-1}Q(f),I_i)\\
\text{(because $f_i\cdot u_{I_{i+1}}=I_i$) }&=\Phi(f_i\cdot (Q(f), I_{i+1})).
\end{align*}
Finally, when $z=f_{r+1}$, we see that (Note that $i\leq r$, so $[e_i,f_{r+1}]=0$):
\begin{align*}
f_{r+1}\cdot\Phi( Q(f),I_{i+1})&=e_if_{r+1}\cdot \Phi( Q(f), I_i)-v^{-1}f_{r+1}\cdot\Phi([e_i,Q(f)], I_{i})\\
\text{(by induction hypothesis)}&=e_i\cdot \Phi(f_{r+1}Q(f), I_i)+v^{\mu_{r+2}-\mu_{r+1}}e_i\cdot \Phi( Q(f), (I_i\cup\{r+2\})\sm\{r+1\})\\
&~-v^{-1}\cdot\Phi([e_i,f_{r+1}Q(f)], I_{i})-v^{\mu_{r+2}-\mu_{r+1}-1+\delta_{i,r}}\Phi([e_i,Q(f)], (I_i\cup\{r+2\})\sm\{r+1\})\\
\text{(by Lemma \ref{lem:Phi-ei})}&=\Phi(f_{r+1}Q(f),I_{i+1}) + v^{\mu_{r+2}-\mu_{r+1}} (1-\delta_{i,r})\Phi(Q(f), (I_{i+1}\cup\{r+2\})\sm\{r+1\}))\\
&=\Phi(f_{r+1}\cdot (Q(f), I_{i+1})).
\end{align*}

\underline{\textbf{Step $2$: $z=\ell$}}

We start the induction step with the equality from \eqref{eq:Phi-I+1}:
\[\Phi( Q(f),I_{i+1})=e_i\cdot \Phi( Q(f), I_i)-v^{-1}\Phi([e_i,Q(f)], I_{i}).\]
We first let $i=1$ in the above equality. Then, we have:
\[
\ell\cdot \Phi( Q(f),I_{2}) = \ell e_1\cdot \Phi( Q(f), I_1)-v^{-1}\ell\cdot \Phi([e_1,Q(f)], I_{1})
\]
We induct on $\ol{\deg}(Q(f))$. When $Q(f)=1$, we have by \eqref{eq:Phi1}: %the LHS becomes:
\begin{align*}\ell\cdot \Phi( 1,I_{2})&=v^{\lambda_1-\lambda_{r+1}}\ell e_1\ell f_1\cdots f_{r-1}f_r\cdot w_{\lambda, r}\\
&=v^{\lambda_1-\lambda_{r+1}}\ell e_1 f_1\cdots f_{r-1}f_r\cdot w_{\lambda, r}\\
&= v^{\lambda_1-\lambda_{r+1}}\ell \left(f_1e_1+\frac{k_1-k_1^{-1}}{v-v^{-1}}\right)f_2\cdots f_{r-1}f_r\cdot w_{\lambda, r},\end{align*}
which is zero since $\ell$ and $e_1$ commute past all the terms $f_2,\ldots, f_r$ and act on $w_{\lambda, r}$ by zero. Next, if $\ol\deg(Q(f))\geq 1$, we can write $Q(f)$ as a linear combination of elements of the form $f_{j}\widetilde Q(f)$ for varying $j$ and $\ol\deg(\widetilde Q(f))< \ol\deg(Q(f)$. When $j\neq 1, r+1$, we have:
\begin{align*}
\ell\cdot \Phi(f_j\widetilde Q(f),I_{2}) &= \ell e_1\cdot \Phi( f_j\widetilde Q(f), I_1)-v^{-1}\ell\cdot \Phi([e_1,f_j\widetilde Q(f)], I_{1})\\
\text{(by \eqref{eq:Phi-fj}) }&=f_j\ell\cdot( e_1\cdot \Phi( \widetilde Q(f), I_1)-v^{-1} \Phi([e_1,\widetilde Q(f)], I_{1}))\\
&=f_j\ell\cdot(\Phi(\widetilde{Q}(f),I_2))\\
\text{(by induction) }&=0.
\end{align*}
Next, suppose $j=r+1$. Then, 
\begin{align*}
\ell\cdot \Phi(f_{r+1}\widetilde Q(f),I_{2}) &= \ell e_1\cdot \Phi( f_{r+1}\widetilde Q(f), I_1)-v^{-1}\ell\cdot \Phi([e_1,f_{r+1}\widetilde Q(f)], I_{1})\\
\text{(by \eqref{eq:Phi-fr}) }&=v^{\mu_{1}-\mu_{r+1}-1}\ell e_1 \ell f_1f_2\cdots f_{r+1}\cdot \Phi( \widetilde Q(f), I_{r+1})-v^{\mu_{1}-\mu_{r+1}-1}\ell\ell f_1f_2\cdots f_{r+1}\cdot \Phi([e_1,\widetilde Q(f)], I_{r+1})\\
&=v^{\mu_{1}-\mu_{r+1}-1}\ell e_1 f_1f_2\cdots f_{r+1}\cdot \Phi( \widetilde Q(f), I_{r+1})-v^{\mu_{1}-\mu_{r+1}-1}\ell f_1f_2\cdots f_{r+1}\cdot \Phi([e_1,\widetilde Q(f)], I_{r+1})\\
&=v^{\mu_{1}-\mu_{r+1}-1}\ell f_1 e_1f_2\cdots f_{r+1}\cdot \Phi( \widetilde Q(f), I_{r+1})-v^{\mu_{1}-\mu_{r+1}-1}\ell f_1f_2\cdots f_{r+1}\cdot \Phi([e_1,\widetilde Q(f)], I_{r+1})\\
&\text{(Here, we use the fact that }\ell\cdot \Phi(\widetilde{Q}(f),I_{r+1})=0 \text{ in order to skip the commutator term}.)\\
&=v^{\mu_{1}-\mu_{r+1}-1}\ell f_1 f_2\cdots f_{r+1}e_1\cdot \Phi( \widetilde Q(f), I_{r+1})-v^{\mu_{1}-\mu_{r+1}-1}\ell f_1f_2\cdots f_{r+1}\cdot \Phi([e_1,\widetilde Q(f)], I_{r+1})\\
\text{(since $e_1\cdot u_{I_{r+1}}=0$)}&=0.
\end{align*}

Lastly, we consider the case when $j=1$. By Corollary~\ref{cor:lastcor} we can assume that $Q(f)$ is of the form $f_1^2Q'(f)$ for some non-commutative monomial $Q'(f)$ (having degree lesser than that of $Q(f)$) or $Q(f)=f_1f_2\cdots f_s$ for some $1\leq s\leq n-1$. In the former case, we have that:
\begin{align*}
\ell\cdot \Phi(f_{1}^2 Q'(f),I_{2}) &= \ell\cdot(\Phi(f_1\cdot(f_1Q'(f),I_2))-v^{\mu_2-\mu_1-2}\Phi(f_1\widetilde{Q}(f),I_1))\\
\text{(by the induction assumption)}&=\ell f_1\cdot\Phi(f_1Q'(f), I_2) - v^{\mu_2-\mu_1-2}\ell\cdot\Phi(f_1Q'(f),I_1)\\
\text{(by the induction assumption)}&= \ell f_1^2\cdot\Phi(Q'(f), I_2) - v^{\mu_2-\mu_1-4}\ell f_1\cdot \Phi(Q'(f),I_1)\\
&- v^{\mu_2-\mu_1-2}\Phi(f_1Q'(f),I_1)\\
\text{(by the induction assumption and $f_1\cdot u_{I_1}=0$)}&= \ell f_1^2\cdot\Phi(Q'(f), I_2) - v^{\mu_2-\mu_1-4}\ell \cdot  \Phi(f_1Q'(f),I_1)\\
&- v^{\mu_2-\mu_1-2}\Phi(f_1Q'(f),I_1)\\
&=\ell f_1^2\cdot\Phi(Q'(f), I_2) - v^{\mu_2-\mu_1-4}(1+v^2)\ell \cdot  \Phi(f_1Q'(f),I_1)\\
\text{(by Equation~\eqref{eq:mun.5})}&=((1+v^2)f_1\ell f_1 - v^2f_1^2\ell)\cdot\Phi(Q'(f), I_2) \\
&- v^{\mu_2-\mu_1-4}(1+v^2)\ell \cdot  \Phi(f_1Q'(f),I_1)\\
\text{(since $\ell\cdot \Phi(Q'(f), I_2)=0$ by induction)}&=(1+v^2)f_1\ell f_1 \cdot\Phi(Q'(f), I_2) - v^{\mu_2-\mu_1-4}(1+v^2)\ell \cdot  \Phi(f_1Q'(f),I_1)\\
&=(1+v^2)f_1\ell \cdot \Phi(f_1Q'(f), I_2) + (1+v^2)v^{\mu_2-\mu_1-4}f_1\ell\cdot \Phi(Q'(f),I_1) \\
&- v^{\mu_2-\mu_1-4}(1+v^2)\ell \cdot  \Phi(f_1Q'(f),I_1)\\
\text{(since $\ell\cdot \Phi(f_1Q'(f), I_2)=0$)}&=0.
\end{align*}

Next, we deal with the case when $Q(f)=f_1f_2\cdots f_s$ for some $s$. If $ s<r+1$, we note that:
\begin{align}
\notag\ell\cdot\Phi(f_1f_2\cdots f_s,I_{2})&=\ell e_1\cdot \Phi(f_1f_2\cdots f_s, I_1)-v^{-1}\ell\cdot \Phi([e_1,f_1f_2\cdots f_s], I_{1})\\
\notag&=\ell e_1\cdot \Phi(f_1f_2\cdots f_s, I_1)-v^{-1}\ell\cdot\Phi\Big(\frac{k_1-k_1^{-1}}{v-v^{-1}}f_2\cdots f_s, I_{1}\Big)\\
\notag&=\ell e_1f_1f_2\cdots f_s\cdot \Phi(1, I_1)-v^{-1}\ell\frac{vk_1-v^{-1}k_1^{-1}}{v-v^{-1}}f_2\cdots f_s\cdot \Phi(1, I_{1})\\
\notag&\text{(by induction assumption)}\\
\notag\text{(by~\eqref{eq:Phi1})}&=v^{\lambda_1-\lambda_{r+1}}\ell e_1f_1f_2\cdots f_s\ell f_1f_2\cdots f_r \cdot  w_{\lambda,r}-v^{\lambda_1-\lambda_{r+1}-1}\ell\frac{vk_1-v^{-1}k_1^{-1}}{v-v^{-1}}f_2\cdots f_s\ell f_1f_2\cdots f_r \cdot  w_{\lambda,r}\\
\notag\text{(by Corollary~\ref{cor:fmove})}&=\frac{v^{\lambda_1-\lambda_{r+1}-1}}{(v+v^{-1})^s}\Big(\ell e_1\ell f_1^2f_2^2\cdots f_s^2 f_{s+1}\cdots f_r -(v+v^{-1})\ell\frac{vk_1-v^{-1}k_1^{-1}}{v-v^{-1}}\ell f_1f_2^2\cdots f_s^2 f_{s+1}\cdots f_r\Big) \cdot  w_{\lambda,r}\\
\notag\text{(by~\eqref{eq:mun.1} and~\eqref{eq:mun.3})}&=\frac{v^{\lambda_1-\lambda_{r+1}-1}}{(v+v^{-1})^s}\ell\Big(e_1 f_1^2 -(v+v^{-1})\frac{vk_1-v^{-1}k_1^{-1}}{v-v^{-1}} f_1\Big)f_2^2\cdots f_s^2 f_{s+1}\cdots f_r \cdot  w_{\lambda,r}.
\end{align}
Evaluating the term inside the middle parentheses in the last line above:
\begin{align*}
 e_1 f_1^2-(v+v^{-1})\frac{vk_1-v^{-1}k_1^{-1}}{v-v^{-1}} f_1& =  f_1e_1f_1+\frac{k_1-k_1^{-1}}{v-v^{-1}}f_1-(v+v^{-1})\frac{vk_1-v^{-1}k_1^{-1}}{v-v^{-1}} f_1\\
& = f_1^2e_1+f_1\frac{k_1-k_1^{-1}}{v-v^{-1}}+ \frac{k_1-k_1^{-1}}{v-v^{-1}}f_1-(v+v^{-1})\frac{vk_1-v^{-1}k_1^{-1}}{v-v^{-1}} f_1\\
& = f_1^2e_1+\frac{v^2k_1-v^{-2}k_1^{-1}}{v-v^{-1}}f_1+\frac{k_1-k_1^{-1}}{v-v^{-1}}f_1-(v+v^{-1})\frac{vk_1-v^{-1}k_1^{-1}}{v-v^{-1}} f_1\\
& = f_1^2e_1.
\end{align*}
Since $e_1$ commutes past all $f_j$'s for $j>1$ to act on $w_{\lambda, r}$ by zero, we conclude that $\ell\cdot \Phi(f_1f_2\cdots f_s, I_2)=0$. Next, we suppose that $s\geq r+1$. In that case, we again have:
\begin{align*}
\ell\cdot\Phi(f_1f_2\cdots f_s,I_{2})&=\ell e_1\cdot \Phi(f_1f_2\cdots f_s, I_1)-v^{-1}\ell\cdot\Phi\Big(\frac{k_1-k_1^{-1}}{v-v^{-1}}f_2\cdots f_s, I_{1}\Big)\\
&=\ell e_1f_1f_2\cdots f_r\cdot \Phi(f_{r+1}f_{r+2}\cdots f_s, I_1)-v^{-1}\ell\frac{vk_1-v^{-1}k_1^{-1}}{v-v^{-1}}f_2\cdots f_r\cdot \Phi(f_{r+1}f_{r+2}\cdots f_s, I_{1})\\
&\text{(by induction assumption)}\\
\text{(by~\eqref{eq:Phi-fr})}&=v^{\lambda_{r+1}-\lambda_1+1}\Big(\ell e_1f_1f_2\cdots f_r-v^{-1}\ell\frac{vk_1-v^{-1}k_1^{-1}}{v-v^{-1}}f_2\cdots f_r\Big)\ell f_1f_2\cdots f_{r+1} \cdot\Phi(f_{r+2}\cdots f_s, I_{r+1})\\
&=v^{\lambda_{r+1}-\lambda_1+1}\Big(\ell e_1f_1f_2\cdots f_r-v^{-1}\ell\frac{vk_1-v^{-1}k_1^{-1}}{v-v^{-1}}f_2\cdots f_r\Big)\ell f_1f_2\cdots f_{r+1}f_{r+2}\cdots f_s \cdot\Phi(1, I_{r+1})\\
&\text{(by induction assumption)}\\
\text{(by~\eqref{eq:topvect})}&=v^{\lambda_{r+1}-\lambda_1+1}\Big(\ell e_1f_1f_2\cdots f_r-v^{-1}\ell\frac{vk_1-v^{-1}k_1^{-1}}{v-v^{-1}}f_2\cdots f_r\Big)\ell f_1f_2\cdots f_{r+1}f_{r+2}\cdots f_s \cdot w_{\lambda, r},
\end{align*}
which is zero by exactly the same argument as in the case when $s<r+1$ using Corollary~\ref{cor:fmove}. This completes the proof of the fact that:
\begin{equation} \label{eq:ell20}\ell\cdot \Phi(Q(f), I_2) = 0 = \Phi(\ell\cdot (Q(f), I_2)).\end{equation}
Finally, we suppose that $i>1$ (so $e_i$ commutes with $\ell$). Then, we have by equation~\eqref{eq:Phi-I+1}:
\[\ell \cdot \Phi(Q(f),I_{i+1})=\ell e_i\cdot \Phi( Q(f), I_i)-v^{-1}\ell\cdot \Phi([e_i,Q(f)], I_{i})=0,\]
where the last equality follows by the induction assumption.

\underline{\textbf{Step $3$: $z=e_j$}}

Let $z=e_j$ where $j\not\in\{i-1,i,i+1\}$. Then, 
\begin{align*}\Phi( e_j\cdot(Q(f),I_{i+1}))=v^{\delta_{j,r+1}}\Phi( e_jQ(f),I_{i+1})&=v^{\delta_{j,r+1}}(e_i\cdot \Phi( e_jQ(f), I_i)-v^{-1}\Phi([e_i,e_jQ(f)], I_{i}))\\
\text{(because $e_j\cdot u_{I_i}=0$ and $[e_i,e_j]=0$) }&=e_i\cdot \Phi(e_j\cdot(Q(f), I_i))-v^{\delta_{j,r+1}-1}\Phi(e_j[e_i,Q(f)], I_{i})\\
\text{(by the induction step and because $e_j\cdot u_{I_i}=0$) }&=e_ie_j\cdot \Phi(Q(f), I_i)-v^{-1}\Phi(e_j\cdot([e_i,Q(f)], I_{i}))\\
\text{(by the induction step and because $[e_i,e_j]=0$) }&=e_je_i\cdot \Phi(Q(f), I_i)-v^{-1}e_j\cdot\Phi([e_i,Q(f)], I_{i}) \\
&=  e_j\cdot\Phi(Q(f),I_{i+1}).\end{align*}

Next, we first consider the case when $i\geq 2$ and $z=e_{i-1}$.
\begin{align*}
&~(v+v^{-1})e_{i-1}\cdot \Phi( Q(f),I_{i+1})\\
=&~(v+v^{-1})e_{i-1}e_i\cdot \Phi( Q(f), I_i)-v^{-1}(v+v^{-1})e_{i-1}\cdot\Phi([e_i,Q(f)], I_{i})\\
=&~(v+v^{-1})e_{i-1}e_i\cdot(e_{i-1}\cdot \Phi( Q(f), I_{i-1})-v^{-1}\Phi([e_{i-1},Q(f)], I_{i-1})) -(1+v^{-2})e_{i-1}\cdot\Phi([e_i,Q(f)], I_{i})\\
=&~(v+v^{-1})e_{i-1}e_ie_{i-1}\cdot \Phi( Q(f), I_{i-1})-(1+v^{-2})e_{i-1}e_i\cdot\Phi([e_{i-1},Q(f)], I_{i-1}) \\
&-(1+v^{-2})e_{i-1}\cdot\Phi([e_i,Q(f)], I_{i})\\
=&~e_{i-1}^2e_i\cdot \Phi( Q(f), I_{i-1})+e_ie_{i-1}^2\cdot \Phi( Q(f), I_{i-1})-(1+v^{-2})e_{i-1}e_i\cdot\Phi([e_{i-1},Q(f)], I_{i-1})\\
&-(1+v^{-2})e_{i-1}\cdot\Phi([e_i,Q(f)], I_{i})\\
=&~e_{i-1}^2\cdot \Phi([e_i,Q(f)], I_{i-1})+e_ie_{i-1}^2\cdot \Phi( Q(f), I_{i-1})-(1+v^{-2})e_{i-1}\cdot\Phi([e_i,[e_{i-1},Q(f)]], I_{i-1})\\
&-(v+v^{-1})\Phi([e_{i-1},[e_i,Q(f)]], I_{i})\\
=&~v^{-2}\Phi([e_{i-1}^2,[e_i,Q(f)]], I_{i-1})+(v+v^{-1})\Phi([e_{i-1},[e_i,Q(f)]], I_i)+v^{-2}e_i\cdot \Phi([e_{i-1}^2,Q(f)], I_{i-1})\\
&+(v+v^{-1})e_i\cdot \Phi([e_{i-1},Q(f)], I_{i})-(1+v^{-2})e_{i-1}\cdot\Phi([e_i,[e_{i-1},Q(f)]], I_{i-1}) -(v+v^{-1})\Phi([e_{i-1},[e_i,Q(f)]], I_{i})\\
=&~v^{-2}\Phi([e_{i-1}^2,[e_i,Q(f)]], I_{i-1})+(v+v^{-1})\Phi([e_{i-1},[e_i,Q(f)]], I_i)+v^{-2} \Phi([e_i,[e_{i-1}^2,Q(f)]], I_{i-1})\\
&+(1+v^{-2})\Phi([e_i,[e_{i-1},Q(f)]], I_{i})+(v+v^{-1})\Phi([e_{i-1},Q(f)], I_{i+1})-v^{-1}(1+v^{-2})\Phi([e_{i-1},[e_i,[e_{i-1},Q(f)]]], I_{i-1})\\
&-(1+v^{-2})\Phi([e_i,[e_{i-1},Q(f)]], I_{i})-(v+v^{-1})\Phi([e_{i-1},[e_i,Q(f)]], I_{i})\\
=&~(v+v^{-1})\Phi([e_{i-1},Q(f)], I_{i+1})\\
=&~(v+v^{-1})\Phi(e_{i-1}\cdot (Q(f), I_{i+1})).
\end{align*}
The case when $z=e_i$ is similar and has been omitted. (When $i=1$, the computation uses relation~\eqref{eq:mun.4} instead of~\eqref{eq:usln.5} as well as equation~\eqref{eq:ell20}.) %Check email on May 6th for precise computation details.

Finally, when $z=e_{i+1}$ and $i\leq r-1$, we have the equality:
\[e_{i+1}\cdot\Phi( Q(f),I_{i+1}) = \Phi(Q(f),I_{i+2})+v^{-1}\Phi([e_{i+1},Q(f)], I_{i+1}),\]
by \eqref{eq:Phi-I+1} with index $i+1$ instead of $i$, proving the required claim. When $i=r$, notice that $e_{r+1}\cdot (Q(f), I_{r+1}) = ([e_{r+1},Q(f)], I_{r+1})$, so we need to prove the equality:
\begin{equation}\label{eq:er+1}e_{r+1}\cdot \Phi(Q(f), I_{r+1}) = \Phi([e_{r+1},Q(f)], I_{r+1}).\end{equation}
When $Q(f)=1$, the RHS of \eqref{eq:er+1} is zero, whereas the LHS is equal to:
\[e_{r+1}\cdot \Phi(1, I_{r+1})= e_{r+1}\cdot w_{\lambda,r}=0,\]
where the first equality follows from~\eqref{eq:topvect} and the last equality from the fact that $w_{\lambda,r}$ is a highest weight vector. %follows by the construction of $V_{\lambda, r}$. 
Next, we suppose that $\ol\deg(Q(f))\geq 1$. We can write $Q(f)$ as a linear combination of elements of the form $f_{j'}\widetilde Q(f)$ for varying $j'$ and $\ol\deg(\widetilde Q(f))< \ol\deg(Q(f))$. Replacing $Q(f)$ by $f_{j'}\widetilde Q(f)$ in \eqref{eq:er+1}, it is clear that the equality follows by the induction assumption as long as $j'\neq r,r+1$. Suppose $j'=r$. Then,
\begin{align*}
e_{r+1}\cdot \Phi(f_{r}\widetilde Q(f), I_{r+1})&=e_{r+1}f_{r}\cdot \Phi(\widetilde Q(f), I_{r+1}) - v^{\mu_{r+1}-\mu_r-2}e_{r+1}\cdot\Phi(\widetilde{Q}(f),I_r)\\
\text{(because $[e_{r+1},f_r]=0$ and by induction) }&=f_{r}e_{r+1}\cdot \Phi(\widetilde Q(f), I_{r+1}) - v^{\mu_{r+1}-\mu_r-1}\Phi([e_{r+1,}\widetilde{Q}(f)],I_r)\\
\text{(by induction) }&=f_{r}\cdot \Phi([e_{r+1},\widetilde Q(f)], I_{r+1}) - v^{\mu_{r+1}-\mu_r-1}\Phi([e_{r+1,}\widetilde{Q}(f)],I_r)\\
\text{(by induction and by the action of $f_r$) }&=\Phi(f_r[e_{r+1},\widetilde Q(f)], I_{r+1})\\
\text{(because $[e_{r+1},f_r]=0$) }&=\Phi([e_{r+1},f_r\widetilde Q(f)], I_{r+1}),
\end{align*}
which is the required equality. Finally, when $j'=r+1$, we have:
\begin{align*}
e_{r+1}\cdot \Phi(f_{r+1}\widetilde Q(f), I_{r+1})&=e_{r+1}f_{r+1}\cdot \Phi(\widetilde Q(f), I_{r+1})\\
&=f_{r+1}e_{r+1}\cdot \Phi(\widetilde Q(f), I_{r+1}) + \frac{k_{r+1}-k_{r+1}^{-1}}{v-v^{-1}}\cdot \Phi(\widetilde Q(f), I_{r+1})\\
\text{(by induction) }&=f_{r+1}\cdot \Phi([e_{r+1},\widetilde Q(f)], I_{r+1}) + \frac{k_{r+1}-k_{r+1}^{-1}}{v-v^{-1}}\cdot \Phi(\widetilde Q(f), I_{r+1})\\
\text{(by induction and by the actions of $f_{r+1}$ and $k_{r+1}$) }&=\Phi(f_{r+1}[e_{r+1},\widetilde Q(f)], I_{r+1}) + \Phi\Big(\frac{k_{r+1}-k_{r+1}^{-1}}{v-v^{-1}}\widetilde Q(f), I_{r+1}\Big)\\
&=\Phi([e_{r+1},f_{r+1}\widetilde Q(f)], I_{r+1}).
\end{align*}
This proves equivariance of the map $\Phi$ with respect to all the generators, thus proving the lemma.
\end{proof}

\end{document}
